% TOP_PROBLEM: {"id":30007519,"problem_number":"LOCAL-30007519","title":"International business-cycle transmission","category_id":21,"category":{"id":21,"name":"economics","display_name":"Economics","description":"Open economic theory statements, published research agendas, and proposed scoped specifications; question provenance and historical priority are qualified per record.","slug":"economics","order_index":21,"created_at":"2026-10-08T00:00:00Z"},"status":"provenance_pending","external_url":null,"legacy_ids":[],"published":false,"statement":"Identify the contribution of international intermediate-input links to cross-country output correlation under an explicit substitution and financial-policy counterfactual.","economics":{"original_id":8,"field":"Business cycles and macro dynamics","kind":"identification","formalization_basis":"proposed_specification","source_status":"not_verified","priority_status":"not_established","priority_note":"No explicit primary open-problem statement matching this catalogue item was verified; supporting research is not evidence of first formulation.","earliest_verified_reference":null,"current_reference":{"authors":["Julian di Giovanni","Şebnem Kalemli-Özcan","Alvaro Silva","Muhammed A. Yıldırım"],"title":"Pandemic-Era Inflation Drivers and Global Spillovers","year":2023,"url":"https://www.newyorkfed.org/medialibrary/media/research/staff_reports/sr1080.pdf","doi":"10.59576/sr.1080","locator":"Abstract and Section 1, pp. 1–3; July 2024 revision","evidence_summary":"The paper quantifies pandemic inflation and international spillovers in a specified multicountry production-network model. It supplies a partial answer and framework, not an explicit general open-problem statement about business-cycle synchronisation or inflation decomposition."},"formalization_note":"The original item was broad; this fixes a production-network counterfactual. The cited inflation model supports the framework but does not explicitly list this synchronisation problem as open."},"statusQualification":"Provenance pending: an explicit published statement of this exact open question was not verified. This is a catalogue-authored candidate specification, not a certified open conjecture. The original item was broad; this fixes a production-network counterfactual. The cited inflation model supports the framework but does not explicitly list this synchronisation problem as open.","statusReviewedAt":"2026-10-08"}
% FIRST_POSTED: 2026-10-08
% ENTRY_KIND: statement_only
% !TeX program = lualatex
\documentclass[11pt]{article}
\usepackage{fontspec}
\usepackage[a4paper,margin=25mm]{geometry}
\usepackage{amsmath,amssymb,mathtools}
\usepackage{xurl}
\usepackage[hidelinks]{hyperref}
\newcommand{\sref}[2]{\hyperlink{source:#1:#2}{[S#2]}}
\setlength{\emergencystretch}{3em}
\begin{document}
% problemId:30007519
\section*{International business-cycle transmission}\label{30007519}
\subsection{Definitions and mathematical statement}
Fix countries \(c=1,\ldots,C\) with measured sectoral input–output links, bilateral trade costs, cross-border assets, and country-specific policy rules. Let \(\Omega\) denote the global production-input matrix, partitioned into domestic and international blocks. A structural model \(M\) maps sectoral and financial innovations into country log outputs \(y_c\). Construct \(\Omega^0\) by replacing imported intermediate inputs with specified domestic substitutes, holding initial final-demand shares and aggregate financial shocks fixed; all resulting equilibrium prices and quantities must adjust within the model.
For horizon \(H\), define the network contribution to cycle synchronisation between countries \(c\) and \(d\) as
\[
\Delta_{cd}(H)=\operatorname{Corr}_M(\Delta_H y_c,\Delta_H y_d;\Omega)
-\operatorname{Corr}_M(\Delta_H y_c,\Delta_H y_d;\Omega^0),
\]
where \(\Delta_H y_c=y_{c,t+H}-y_{c,t}\), both correlations use the same declared innovation distribution, and each output change has strictly positive finite variance under both networks. Identify \(\Delta_{cd}(H)\), or a sensitivity set, using shocks with independently supported geographic incidence. Distinguish propagation from common shocks and endogenous network formation; specify whether national policy rules remain fixed or optimally respond. Validate predicted transmission against held-out trade disruptions. Removing international links is a model-based intervention whose substitute technology must be disclosed; its correlation effect cannot be read directly from observed trade–cycle correlations.

\par\textbf{Formulation and scope.} The original item was broad; this fixes a production-network counterfactual. The cited inflation model supports the framework but does not explicitly list this synchronisation problem as open.

\par\textbf{Open-question provenance.} An explicit published open-question statement was not verified. Sources below are supporting literature and must not be treated as a first listing.

\par\textbf{Historical priority.} No explicit primary open-problem statement matching this catalogue item was verified; supporting research is not evidence of first formulation.

\subsection{Short English statement}
Identify the contribution of international intermediate-input links to cross-country output correlation under an explicit substitution and financial-policy counterfactual.

\subsection{Sources}
\begin{itemize}\raggedright
\item[\textnormal{[S1]}]\hypertarget{source:30007519:1}{} Julian di Giovanni, Şebnem Kalemli-Özcan, Alvaro Silva, Muhammed A. Yıldırım. 2023. Pandemic-Era Inflation Drivers and Global Spillovers. Abstract and Section 1, pp. 1–3; July 2024 revision.
\url{https://www.newyorkfed.org/medialibrary/media/research/staff_reports/sr1080.pdf}.
DOI: 10.59576/sr.1080.
{\small\textit{Supporting or current literature; see evidence qualification. The paper quantifies pandemic inflation and international spillovers in a specified multicountry production-network model. It supplies a partial answer and framework, not an explicit general open-problem statement about business-cycle synchronisation or inflation decomposition.}}
\end{itemize}
\end{document}
